\chapter{Formação contínua e certificações}
\label{cap:formacao}

\begin{objetivos}
Ao final deste capítulo você deve ser capaz de:
\begin{itemize}
  \item situar sua formação atual em uma trilha de estudo coerente;
  \item distinguir certificação de produto, de competência e de treinamento pontual;
  \item escolher recursos de estudo por critério, e não por propaganda;
  \item montar uma bancada de prática com custo controlado e segurança;
  \item executar um plano de 90 dias com entrega verificável ao final de cada etapa.
\end{itemize}
\end{objetivos}

\section{Por que a formação não termina no curso}

O controlador mudou três vezes em vinte anos: passou de ilha de lógica a nó de rede, de rede a
sistema com camada de informação, e agora convive com visão computacional, borda e modelos de
linguagem (Capítulo~\ref{cap:ti-to}). O profissional que parou em ``programa em Ladder'' fica preso a
uma geração de equipamento.

Duas consequências práticas para quem está se formando:

\begin{enumerate}
  \item \textbf{a base não envelhece}: eletricidade, comandos, lógica, instrumentação e a norma
        IEC 61131-3 --- com suas edições --- sustentam qualquer marca. É onde vale investir primeiro;
  \item \textbf{o entorno muda}: protocolos, ferramentas e plataformas de dados se renovam. Aqui o
        esforço é de atualização contínua, em dose pequena e frequente, e não de recomeço.
\end{enumerate}

\index{formação continuada}

\section{A trilha}

\begin{figure}[H]
  \centering
  \includegraphics[width=0.41\linewidth]{25-trilha-formacao.pdf}
  \caption{Trilha sugerida de estudo, da base elétrica à integração de sistemas.}
  \label{fig:trilha-formacao}
\end{figure}

A Figura~\ref{fig:trilha-formacao} apresenta a sequência; a Tabela~\ref{tab:trilha} detalha cada
etapa e a entrega que comprova o domínio.

\begin{table}[H]
  \centering
  \caption{Etapas da trilha, o que se aprende e como provar que aprendeu.}
  \label{tab:trilha}
  \begin{tabular}{p{3.4cm} p{5.2cm} p{5.7cm}}
    \toprule
    Etapa & Conteúdo & Evidência de domínio \\
    \midrule
    1. Eletricidade e comandos elétricos
      & Leitura de diagrama, contator, relé, selo, botoeira, proteção de circuito
      & Montar um comando de partida e parada com selo em bancada, e explicar cada contato \\
    \addlinespace
    2. Lógica digital e instrumentação
      & Tabela-verdade, álgebra booleana, sinal $4$--$20\,$mA, escalas, precisão
      & Converter um sinal de $4$--$20\,$mA em grandeza de engenharia e justificar a resolução obtida
        (Capítulo~\ref{cap:es}) \\
    \addlinespace
    3. CLP conforme IEC 61131-3
      & Modelo de software, tipos de dados, LD, FBD, ST e SFC, tarefas
      & Escrever o mesmo controle em LD e em ST, com tipos declarados, e documentar a diferença \\
    \addlinespace
    4. Redes industriais
      & Modbus RTU e TCP, protocolos industriais, arquiteturas de rede
      & Ligar dois equipamentos por Modbus e construir a tabela de registros usados
        (Capítulo~\ref{cap:redes-scada}) \\
    \addlinespace
    5. Supervisão e SCADA
      & Telas, tags, histórico, alarmes e eventos
      & Montar uma tela com histórico e um alarme classificado corretamente
        (Capítulo~\ref{cap:redes-scada}) \\
    \addlinespace
    6. Integração, dados e segurança
      & OPC UA, MQTT, indicadores, zonas e condutos
      & Publicar uma variável por protocolo aberto, calcular um indicador e desenhar as zonas da
        bancada (Capítulos~\ref{cap:iiot} e~\ref{cap:seguranca}) \\
    \bottomrule
  \end{tabular}
\end{table}

O erro mais comum na trilha é começar pela etapa 5, porque a tela é a parte visível. Sem a etapa 3,
a tela mostra dado que ninguém sabe se está certo, como a Tabela~\ref{tab:trilha} deixa explícito ao
exigir a entrega de cada etapa antes de seguir.

\section{Certificações, treinamentos e o que eles valem}

Vale separar três coisas que costumam aparecer misturadas em conversa:

\begin{itemize}
  \item \textbf{treinamento de fabricante}: curso sobre uma família de produtos. Ensina profundamente
        o equipamento e é o que o mercado de manutenção procura. Programas de treinamento existem em
        fabricantes de peso no país, como Siemens (SITRAIN), Rockwell Automation, Schneider Electric,
        WEG, Novus, Mitsubishi Electric, Omron e Beckhoff. O conteúdo, a duração e a validade mudam
        com o tempo --- confirme sempre na fonte, no site oficial do fabricante;
  \item \textbf{certificação de competência profissional}: avalia o profissional, não o produto.
        Existem programas internacionais com esse perfil, como os da ISA (\emph{International Society
        of Automation}) voltados a automação e a sistemas de controle. Exigem experiência comprovada
        e são mantidos por recertificação periódica;
  \item \textbf{certificação de produto ou de ferramenta}: atesta que um produto ou ambiente de
        programação atende a requisitos de uma norma ou de uma biblioteca padronizada. É o caso das
        certificações ligadas à PLCopen para ambientes de programação (Capítulo~\ref{cap:familia-iec}).
        \emph{Ela não certifica pessoas} --- não confunda com certificação profissional.
\end{itemize}

\begin{atencao}
Preço, carga horária, pré-requisitos e prazo de validade de certificações mudam com frequência e
variam por país. Nenhuma informação deste capítulo deve ser usada como referência de custo ou de
validade: consulte sempre a entidade emissora antes de se inscrever ou de anunciar um título no
currículo.
\end{atencao}

Três critérios para escolher onde investir tempo e dinheiro:

\begin{enumerate}
  \item \textbf{a instituição tem laboratório?} Curso de automação sem bancada é curso de leitura;
  \item \textbf{o conteúdo ensina norma ou só uma marca?} Um bom curso diz, em cada instrução, se ela
        é normatizada ou proprietária --- é a distinção que o Capítulo~\ref{cap:iec-versus-legado}
        ensina;
  \item \textbf{existe avaliação prática?} Certificado que se obtém apenas respondendo a um
        questionário prova menos.
\end{enumerate}

\section{Recursos de estudo}

A Tabela~\ref{tab:recursos-estudo} reúne o que costuma ser útil e o cuidado correspondente a cada
recurso.

\begin{table}[H]
  \centering
  \caption{Recursos úteis e o cuidado correspondente.}
  \label{tab:recursos-estudo}
  \begin{tabular}{p{3.6cm} p{5.2cm} p{5.5cm}}
    \toprule
    Recurso & Para que serve & Cuidado \\
    \midrule
    Ambientes de programação em versão de avaliação ou de uso educacional
      & Praticar a linguagem e a estrutura de POUs sem comprar licença
      & Nem todo recurso fica disponível na versão gratuita; verifique os limites antes de investir
        tempo em material didático de terceiros \\
    \addlinespace
    Ambientes de desenvolvimento conforme a norma, com versão de uso livre
      & Aprender o modelo de software da norma, independente de fabricante
      & Verifique em qual edição da norma o ambiente se declara conforme \\
    \addlinespace
    Simuladores de processo
      & Testar lógica sem hardware, inclusive casos de falha
      & Simulação não substitui bancada para aprender elétrica, aterramento e ruído \\
    \addlinespace
    Plataformas abertas de controlador (por exemplo, projetos abertos de CLP executando em Linux)
      & Entender o que há dentro de um controlador: varredura, memória, E/S
      & Não são substitutos de controlador industrial certificado em aplicação real \\
    \addlinespace
    Documentação técnica do fabricante
      & A fonte primária: manual de instruções, de comunicação e de diagnóstico
      & É a referência a citar; fórum é apoio, não fonte \\
    \addlinespace
    Normas
      & A definição correta de conceito, linguagem e tipo de dado
      & São documentos pagos, obtidos junto aos organismos de normalização. No Brasil, a ABNT publica
        adoções de normas IEC com a designação NBR IEC --- confira no catálogo quais partes existem e
        em que edição. Nunca use cópia não autorizada \\
    \addlinespace
    Comunidades técnicas e fóruns
      & Resolver problema concreto e ver como outros resolvem
      & Verifique a resposta na documentação; a resposta com mais votos não é necessariamente a
        correta \\
    \bottomrule
  \end{tabular}
\end{table}

Nunca reproduza trechos de norma, figura de catálogo ou manual de fabricante em trabalho publicado
sem autorização. Cite a referência: obra, fabricante, título, versão e data.

\index{normas!ABNT NBR IEC}

\section{Bancada própria}

Uma bancada simples cobre quase todo o conteúdo deste material e cabe em uma mesa:

\begin{itemize}
  \item um controlador compacto, mesmo de porte pequeno, com entrada digital, saída a relé e uma
        entrada analógica;
  \item fonte de alimentação, disjuntor e proteção na entrada --- e a bancada \emph{não} deve ser
        ligada à rede elétrica sem supervisão de quem conhece o risco;
  \item botoeiras, chaves seletoras, sinaleiros, fim de curso e um relé como carga;
  \item um sensor qualquer para a entrada analógica: potenciômetro para simular, transmissor para
        aprender escala;
  \item um computador com o ambiente de programação e a planilha de documentação;
  \item itens de segurança: protetor de bancada, destino claro para os cabos e \textbf{regra escrita
        de que não se monta nem se mede circuito energizado}.
\end{itemize}

Duas adições baratas aumentam muito o valor didático: um módulo de entrada/saída remota acessível por
Modbus, para praticar rede, e uma tela de supervisão simples, para praticar tag, alarme e histórico.

\section{Plano de 90 dias}

O plano da Tabela~\ref{tab:plano-90-dias} é uma sugestão de sequência para quem já fez o básico de
eletricidade. O que importa
não é o prazo, é a regra: \textbf{toda etapa termina com algo que outra pessoa pode ver} --- e não
com horas de vídeo assistido.

\begin{table}[H]
  \centering
  \caption{Plano de estudo de 90 dias, com entrega verificável por etapa.}
  \label{tab:plano-90-dias}
  \begin{tabular}{p{1.8cm} p{4.2cm} p{3.8cm} p{4.5cm}}
    \toprule
    Período & Foco & Prática & Entrega \\
    \midrule
    Semanas 1--2
      & Tipos de dados, endereçamento e leitura de diagrama de comando
      & Montar um comando de partida e parada com selo em bancada
      & Diagrama desenhado, lista dos pontos usados e vídeo curto do funcionamento \\
    \addlinespace
    Semanas 3--4
      & Ladder e texto estruturado conforme a norma
      & Escrever o mesmo controle nas duas linguagens, com tipos declarados
      & Os dois programas versionados, com comentário explicando a diferença \\
    \addlinespace
    Semanas 5--6
      & Temporização, contagem, borda e retenção
      & Implementar uma sequência com etapas, intertravamento e retomada após falta de energia
      & Lista de casos de teste preenchida, com resultado observado \\
    \addlinespace
    Semanas 7--8
      & Rede industrial e supervisão
      & Ligar remota por Modbus, criar tags e montar tela com histórico e alarme
      & Tabela de registros, tela funcionando e alarme classificado \\
    \addlinespace
    Semanas 9--10
      & Documentação e boas práticas
      & Documentar o projeto das semanas anteriores como entrega de engenharia
      & Descritivo, lista de E/S, tabela de tags e rotina de backup preenchida \\
    \addlinespace
    Semanas 11--12
      & Integração e dados
      & Publicar uma variável por OPC UA ou MQTT e calcular um indicador
      & Diagrama da arquitetura, indicador calculado e desenho das zonas da bancada \\
    \bottomrule
  \end{tabular}
\end{table}

Ao final do plano, o conjunto das entregas forma um \textbf{portfólio}. Em processo seletivo, ele vale
mais que uma lista de cursos: mostra o que a pessoa sabe fazer, como documenta e como testa.

\begin{pratica}
Monte o seu plano de 90 dias:
(a) escreva as seis etapas adaptadas à sua realidade, com o tempo que você tem por semana;
(b) defina, para cada etapa, a entrega verificável;
(c) liste o que falta comprar ou conseguir emprestado na bancada, com prioridade;
(d) agende a revisão da etapa 1 para uma data específica --- com alguém que possa assistir à
demonstração.
\end{pratica}

\begin{exercicios}
\begin{enumerate}[label=\alph*)]
  \item Explique por que a base da trilha (etapas 1 a 3) não envelhece, e o que muda nas etapas 4 a 6.
  \item Por que começar pela supervisão é um erro comum, e o que ele produz?
  \item Diferencie treinamento de fabricante, certificação de competência e certificação de produto.
  \item Por que a certificação de um ambiente de programação \emph{não} certifica o profissional?
  \item Cite três critérios para escolher um curso de automação e explique cada um.
  \item Escolha uma das etapas da Tabela~\ref{tab:trilha} e descreva uma entrega sua, real, que sirva
        como evidência de domínio dela.
  \item Por que simulação não substitui bancada no aprendizado de elétrica?
  \item Que cuidado se deve ter ao usar resposta de fórum como referência técnica?
  \item Liste cinco itens mínimos de uma bancada de estudo e um item de segurança obrigatório.
  \item Por que a bancada deve ter regra escrita sobre trabalho em circuito energizado?
  \item Cite duas adições baratas à bancada e diga o que cada uma ensina.
  \item Explique a regra ``toda etapa termina com algo que outra pessoa pode ver'' e por que ela
        reduz o risco de estudo sem prática.
  \item Monte a entrega da semana 5--6 do plano de 90 dias para um processo simples, no formato da
        Tabela~\ref{tab:roteiro-teste}.
  \item Por que um portfólio de entregas pode valer mais que uma lista de cursos em um processo
        seletivo?
\end{enumerate}
\end{exercicios}
